% SGA 6, Exposición VII: fragmento español alineado con la fuente.
% Alcance: sección 2 completa, líneas 9965--10143 inclusive.
% Fuente francesa: sga6_fr_r1/src/fr.tex.
% SHA-256 de la fuente francesa completa:
% 0D2158B9FE3B466309A92165727A0B191422E05FBAA0DACD5806E8322C54BA1A.
\providecommand{\sisi}[2]{#2}

\section*{2. Cálculos sobre las inmersiones regulares}

\subsection*{2.1}
Sean \(A\) un anillo, y \(B\) y \(C\) dos \(A\)-álgebras. Nos proponemos dotar a \(\bigoplus_{i\ge0}\Tor_i^A(B,C)\) de una estructura de \(A\)-álgebra graduada, es decir, definir homomorfismos
\begin{equation}\tag{2.1.1}
 \Tor_i^A(B,C)\otimes_A\Tor_j^A(B,C)\longrightarrow \Tor_{i+j}^A(B,C).
\end{equation}
Sea \(P_\bullet\) una resolución de \(B\) por \(A\)-módulos proyectivos; los \(\Tor_i^A(B,C)\) son los grupos de homología del complejo \(P_\bullet\otimes_A C\). Se sabe definir homomorfismos canónicos ([1], IV 6):
\begin{equation}\tag{2.1.2}
 H_i(P_\bullet\otimes_A C)\otimes_A H_j(P_\bullet\otimes_A C)
 \longrightarrow H_{i+j}\bigl((P_\bullet\otimes_A C)\otimes_A(P_\bullet\otimes_A C)\bigr).
\end{equation}
Por otra parte, \(P_\bullet\otimes_A P_\bullet\) es un complejo de \(A\)-módulos proyectivos, en grados \(\geq0\), y \(H_0(P_\bullet\otimes_A P_\bullet)=B\otimes_A B\); además, \(P_\bullet\) es una resolución de \(B\). Por consiguiente, el morfismo de aumento \(B\otimes_A B\to B\) se prolonga a un morfismo de complejos \(P_\bullet\otimes_A P_\bullet\to P_\bullet\), y define, teniendo en cuenta el homomorfismo \(C\otimes_A C\to C\) deducido de la estructura de álgebra de \(C\), homomorfismos
\begin{equation}\tag{2.1.3}
 H_{i+j}\bigl((P_\bullet\otimes_A C)\otimes_A(P_\bullet\otimes_A C)\bigr)
 \longrightarrow H_{i+j}(P_\bullet\otimes_A C).
\end{equation}
Componiendo (2.1.2) y (2.1.3), se obtienen los homomorfismos (2.1.1). La estructura de álgebra definida por los homomorfismos (2.1.1) es asociativa (cf. [1] XI 2). Por otra parte, el \(A\)-módulo \(H_i(P_\bullet\otimes_A C)\) está dotado de manera estándar de una estructura de \(B\otimes_A C\)-módulo, haciendo operar \(C\) de manera evidente y \(B\) por funtorialidad; de igual modo, los \(H_j\bigl((P_\bullet\otimes_A C)\otimes_A(P_\bullet\otimes_A C)\bigr)\) están dotados de dos estructuras de \(B\otimes_A C\)-módulos. Está claro que los homomorfismos (2.1.2) y (2.1.3) son compatibles con estas dos estructuras: para (2.1.2), por funtorialidad (situándonos en la categoría de los \(A\)-módulos), y para (2.1.3), utilizando el resultado clásico sobre los morfismos de un complejo acíclico. Hemos dotado así a \(\bigoplus_{i\ge0}\Tor_i^A(B,C)\) de una estructura de \(B\otimes_A C\)-álgebra asociativa.

\subsection*{2.2}
Supongamos que \(B\) admite una resolución por \(A\)-módulos proyectivos \(P_\bullet\), tal que \(P_\bullet=\bigoplus_{i\ge0}P_i\) sea una \(A\)-álgebra graduada, satisfaciendo la diferencial:
\[
\begin{array}{ll}
\text{a)}& P_0\longrightarrow B\text{ es un homomorfismo de anillos~;}\\[0.4ex]
\text{b)}& \forall x\in P_i,\ \forall y\in P_j,\quad d(xy)=d(x)y+(-1)^i x\,d(y).
\end{array}
\]
\sgastandalonefootnotetext{x-tor-resolution}{Se puede demostrar que tal resolución existe siempre.}
Entonces las aplicaciones «~producto~» definen diagramas conmutativos para \(i,j\ge0\):
\begin{equation}\tag{2.2.1}
\begin{tikzcd}[column sep=large,row sep=large]
P_i\otimes_A P_j \arrow[r] \arrow[d] & (P_{i-1}\otimes_A P_j)\oplus(P_i\otimes_A P_{j-1}) \arrow[d]\\
P_{i+j} \arrow[r] & P_{i+j-1},
\end{tikzcd}
\end{equation}
y por tanto un homomorfismo de complejos \(P_\bullet\otimes_A P_\bullet\to P_\bullet\) que prolonga el homomorfismo de aumento \(B\otimes_A B\to B\). Esto muestra que la estructura de álgebra de \(\bigoplus_{i\ge0}\Tor_i^A(B,C)\) puede calcularse a partir de la de \(P_\bullet\). En particular, si \(P_\bullet\) es una \(A\)-álgebra alternada (es decir, tal que para todo \(x\in P_\bullet\) se tenga \(x^2=0\)), lo mismo vale para \(\bigoplus_{i\ge0}\Tor_i^A(B,C)\).

Esta última observación se aplica al caso en que \(B\) es el cociente de \(A\) por un ideal regular, pues se puede tomar entonces como \(P_\bullet\) el complejo de Koszul asociado a un homomorfismo regular que define el ideal de \(B\).

\subsection*{2.3}
La estructura de \(B\otimes_A C\)-álgebra definida en 2.1 sobre \(\bigoplus_{i\ge0}\Tor_i^A(B,C)\) es funtorial en \(B\) y \(C\). Se pueden globalizar entonces las construcciones precedentes: sean \(X\) un esquema, y \(X'\), \(X''\) dos \(X\)-esquemas; se pueden definir localmente sobre \(X'\times_X X''\) los morfismos (2.1.1), que se pegan por funtorialidad y definen sobre \(\bigoplus_{i\ge0}\Tor_i^{\cO_X}(\cO_{X'},\cO_{X''})\) una estructura de \(\cO_{X'\times_X X''}\)-álgebra graduada asociativa.

\textbf{Lema 2.4.} Sean \(X',X''\) dos \(X\)-esquemas, y \(M\) un \(\cO_{X'}\)-módulo plano. Entonces, para todo \(i\ge0\), se tiene un isomorfismo canónico
\begin{equation}\tag{2.4.1}
\Tor_i^{\cO_X}(M,\cO_{X''})\simeq
M\otimes_{\cO_{X'}}\Tor_i^{\cO_X}(\cO_{X'},\cO_{X''}).
\end{equation}
Basta definir un isomorfismo funtorial en el caso afín; se podrá entonces, por pegado, definirlo globalmente en el caso general. Sean, pues, \(X=\Spec(A)\), \(X'=\Spec(A')\), \(X''=\Spec(A'')\); \(M\) está definido por un \(A'\)-módulo plano \(M\). Sea \(P_\bullet\) una resolución proyectiva de \(A''\) sobre \(A\); los \(\Tor_i^A(A',A'')\) son los grupos de homología de \(A'\otimes_A P_\bullet\). Como \(M\) es plano sobre \(A'\), se tiene un isomorfismo canónico \(M\otimes_{A'}\Tor_i^A(A',A'')\xrightarrow{\sim}H_i\bigl(M\otimes_{A'}(A'\otimes_A P_\bullet)\bigr)\). Por otra parte, el isomorfismo canónico \(M\otimes_{A'}(A'\otimes_A P_\bullet)\xrightarrow{\sim}M\otimes_A P_\bullet\) da un isomorfismo canónico \(H_i\bigl(M\otimes_{A'}(A'\otimes_A P_\bullet)\bigr)\xrightarrow{\sim}\Tor_i^A(M,A'')\), de donde resulta la afirmación.

\textbf{Proposición 2.5.} Sean \(X\) un esquema, \(Y\) y \(Z\) dos subesquemas cerrados de \(X\), y \(T=Y\times_X Z\). Se supone que, para todo \(x\in T\), existen un entorno abierto \(U\) de \(x\) y una sucesión de secciones \(f_1,\ldots,f_p,g_1,\ldots,g_r,h_1,\ldots,h_s\) de \(\cO_U\), tales que el complejo de Koszul asociado a \(f_1,\ldots,f_p,g_1,\ldots,g_r\) sea una resolución de \(\cO_{Y\cap U}\), el complejo de Koszul asociado a \(f_1,\ldots,f_p,h_1,\ldots,h_s\) una resolución de \(\cO_{Z\cap U}\), y el complejo de Koszul asociado a \(f_1,\ldots,f_p,g_1,\ldots,g_r,h_1,\ldots,h_s\) una resolución de \(\cO_{T\cap U}\) (condiciones satisfechas, en particular, en el caso noetheriano si la sucesión \(f_1,\ldots,f_p,g_1,\ldots,g_r,h_1,\ldots,h_s\) es regular y si \(Y=V(f_1,\ldots,f_p,g_1,\ldots,g_r)\), \(Z=V(f_1,\ldots,f_p,h_1,\ldots,h_s)\)). Sean \(I,J\) los ideales de \(Y,Z\). Se tienen entonces los isomorfismos siguientes:
\[
\begin{array}{ll}
\text{i)}& \forall i\ge0,\quad
\Tor_i^{\cO_X}(\cO_Y,\cO_Z)\xrightarrow{\sim}
\bigwedge^i\Tor_1^{\cO_X}(\cO_Y,\cO_Z);\\[0.8ex]
\text{ii)}&
\Tor_1^{\cO_X}(\cO_Y,\cO_Z)\simeq
I\cap J/I.J\simeq
\Ker\bigl(i^*(\sisi{N_{Y/X}}{N_{X/Y}})\times j^*(\sisi{N_{Z/X}}{N_{X/Z}})\to \sisi{N_{T/X}}{N_{X/T}}\bigr), % Canon SGA-SRC-001489
\end{array}
\]
donde se denota por \(i\) la inmersión de \(T\) en \(Y\), y por \(j\) la inmersión de \(T\) en \(Z\).

En particular, para toda inmersión cerrada regular \(Y\to X\), se tiene
\[
\begin{array}{ll}
\text{i')}& \forall i\ge0,\quad
\Tor_i^{\cO_X}(\cO_Y,\cO_Y)\xrightarrow{\sim}
\bigwedge^i\Tor_1^{\cO_X}(\cO_Y,\cO_Y);\\[0.8ex]
\text{ii')}&
\Tor_1^{\cO_X}(\cO_Y,\cO_Y)\simeq \sisi{N_{Y/X}}{N_{X/Y}}.
\end{array}
\]

\emph{Demostración.} i) Hemos visto en 2.3 que \(\bigoplus_{i\ge0}\Tor_i^{\cO_X}(\cO_Y,\cO_Z)\) podía dotarse de una estructura de \(\cO_T\)-álgebra; por 2.2, esta álgebra es alternada, pues la inmersión de \(Y\) (o de \(Z\)) en \(X\) es regular. Existe, pues, un único homomorfismo de álgebras
\begin{equation}\tag{2.5.1}
\bigwedge \Tor_1^{\cO_X}(\cO_Y,\cO_Z)\longrightarrow
\bigoplus_{i\ge0}\Tor_i^{\cO_X}(\cO_Y,\cO_Z),
\end{equation}
que prolonga la identidad en grado \(1\).

Para ver que es un isomorfismo, se puede comprobar localmente. Nos situamos, pues, en un entorno de un punto de \(T\), y consideramos secciones \(f_1,\ldots,f_p,g_1,\ldots,g_r,h_1,\ldots,h_s\) de \(\cO_X\) en este entorno, con las propiedades indicadas en el enunciado de 2.5. Pongamos \(E=(\cO_X)^p\), \(E'=(\cO_X)^r\), y sean \(u:E\to\cO_X\) el homomorfismo definido por los \(f_i\), \(v:E'\to\cO_X\) el homomorfismo definido por los \(g_j\), y \(K_\bullet(u)\) y \(K_\bullet(v)\) los complejos de Koszul asociados. El complejo de Koszul asociado a la sucesión \(f_1,\ldots,f_p,g_1,\ldots,g_r\) es isomorfo al complejo producto tensorial \(K_\bullet(u)\otimes_{\cO_X}K_\bullet(v)\). Como es una resolución de \(\cO_Y\) por \(\cO_X\)-módulos libres de tipo finito, los \(\Tor_i^{\cO_X}(\cO_Y,\cO_Z)\) son los haces de homología del complejo \(K_\bullet(u)\otimes_{\cO_X}K_\bullet(v)\otimes_{\cO_X}\cO_Z\). Se puede considerar como el producto tensorial de los complejos \(K_\bullet(u)\) y \(K_\bullet(v)\otimes_{\cO_X}\cO_Z\), y utilizar para calcular su homología las sucesiones espectrales asociadas a un bicomplejo (EGA \(0_{\mathrm{III}}\) (11.3.5)), cuyo límite es la homología del complejo simple asociado.

Consideremos la sucesión espectral tal que \({}'E^2_{pq}=H^I_p\bigl(H^{II}_q(K_{\bullet,\bullet})\bigr)\), con las notaciones de EGA \(0_{\mathrm{III}}\) (11.3.5), siendo el primer índice relativo al complejo \(K_\bullet(u)\) y el segundo al complejo \(K_\bullet(v)\otimes_{\cO_X}\cO_Z\). Ahora bien, el complejo \(K_\bullet(v)\otimes_{\cO_X}\cO_Z\) es el complejo de Koszul sobre \(\cO_Z\) asociado a la sucesión de secciones \(\overline g_1,\ldots,\overline g_r\), donde los \(\overline g_i\) son las imágenes de los \(g_i\) en \(\cO_Z\), y este complejo es una resolución de \(\cO_T\) sobre \(\cO_Z\), pues las inmersiones \(Z\to X\) y \(T\to Z\to X\) son regulares (cf. demostración de 1.7 ii). Como \(\bigwedge^p E\) es libre sobre \(\cO_X\), el complejo \(\bigwedge^pE\otimes_{\cO_X}(K_\bullet(v)\otimes_{\cO_X}\cO_Z)\) es una resolución de \(\bigwedge^pE\otimes_{\cO_X}\cO_T\). Por consiguiente, la sucesión espectral degenera, y los \(\Tor_i^{\cO_X}(\cO_Y,\cO_Z)\) son los haces de homología del complejo \(K_\bullet(u)\otimes_{\cO_X}\cO_T\). Ahora bien, la diferencial de \(K_\bullet(u)\) se expresa en términos de los \(f_i\), que anulan \(\cO_T\); por tanto, la diferencial de \(K_\bullet(u)\otimes_{\cO_X}\cO_T\) es nula, y se obtienen los isomorfismos
\begin{equation}\tag{2.5.2}
\underline{\Tor}_i^{\cO_X}(\cO_Y,\cO_Z)
\xrightarrow{\sim}
\cO_T\otimes_{\cO_X}\bigwedge^iE
\longrightarrow
\bigwedge^i(\cO_T)^p.
\end{equation}
En particular, para \(i=1\), se obtiene
\[
\Tor_1^{\cO_X}(\cO_Y,\cO_Z)\xrightarrow{\sim}\sisi{(\cO_T)^n}{(\cO_T)^p},
\]
de donde, para todo \(i\):
\begin{equation}\tag{2.5.3}
\underline{\Tor}_i^{\cO_X}(\cO_Y,\cO_Z)
\xrightarrow{\sim}
\sisi{\bigwedge^i\Tor_i^{\cO_X}(\cO_Y,\cO_Z)}{\bigwedge^i\Tor_1^{\cO_X}(\cO_Y,\cO_Z)}.
\end{equation}

Ahora bien, resulta de 2.2 que el isomorfismo
\[
\bigwedge(\cO_T)^n\xrightarrow{\sim}
\bigoplus_{i\ge0}\Tor_i^{\cO_X}(\cO_Y,\cO_Z)
\]
definido por los isomorfismos (2.5.2) es compatible con las estructuras de álgebra de ambos miembros; lo mismo vale, pues, para el isomorfismo
\[
\sisi{\bigwedge\Tor_i^{\cO_X}(\cO_Y,\cO_Z)}{\bigwedge\Tor_1^{\cO_X}(\cO_Y,\cO_Z)}\xrightarrow{\sim}
\bigoplus_{i\ge0}\Tor_i^{\cO_X}(\cO_Y,\cO_Z)
\]
definido por los isomorfismos (2.5.3). Por razones de unicidad, este último es la restricción del homomorfismo (2.5.1) que había sido definido globalmente. Por tanto, el homomorfismo (2.5.1) es un isomorfismo.

ii) Consideremos la sucesión exacta
\[
0\longrightarrow I\longrightarrow \cO_X\longrightarrow \cO_Y\longrightarrow0,
\]
y la sucesión exacta de los Tor que resulta de ella:
\[
0\longrightarrow\underline{\Tor}_1^{\cO_X}(\cO_Y,\cO_Z)
\longrightarrow I\otimes_{\cO_X}\cO_Z
\longrightarrow \cO_Z.
\]
Se sigue que \(\underline{\Tor}_1^{\cO_X}(\cO_Y,\cO_Z)\) es el núcleo del homomorfismo
\[
I\otimes_{\cO_X}(\cO_X/J)\longrightarrow \cO_X/J,
\]
es decir, \(I\cap J/I.J\).

Para demostrar la última relación, vamos a demostrar que existe una sucesión exacta
\[
0\longrightarrow I\cap J/I.J
\longrightarrow i^*(\sisi{N_{Y/X}}{N_{X/Y}})\times j^*(\sisi{N_{Z/X}}{N_{X/Z}})
\longrightarrow \sisi{N_{T/X}}{N_{X/T}}\longrightarrow0.
\]
Los homomorfismos \(i^*(\sisi{N_{Y/X}}{N_{X/Y}})\to \sisi{N_{T/X}}{N_{X/T}}\) y \(j^*(\sisi{N_{Z/X}}{N_{X/Z}})\to \sisi{N_{T/X}}{N_{X/T}}\) se definieron en (1.6.1); tomamos aquí su diferencia. Por otra parte, \(I\cap J\) se envía naturalmente en \(I/I^2\), y por tanto en \(i^*(I/I^2)\), y del mismo modo en \(j^*(J/J^2)\), y estos homomorfismos se anulan sobre \(I.J\), pues \(I,J\subset K\) (donde \(K\) designa el ideal de \(T\)); esto define la flecha de la izquierda.

Para comprobar que la sucesión es exacta, se puede hacer localmente sobre \(T\), utilizando las secciones definidas en el enunciado de 2.5. Se tiene entonces una base de \(\sisi{N_{T/X}}{N_{X/T}}\) formada por las imágenes de los \(f_i,g_j,h_k\), una base de \(i^*(\sisi{N_{Y/X}}{N_{X/Y}})\) formada por las imágenes de los \(f_i,g_j\), y una base de \(j^*(\sisi{N_{Z/X}}{N_{X/Z}})\) formada por las imágenes de los \(f_i,h_k\). Está claro entonces que \(i^*(\sisi{N_{Y/X}}{N_{X/Y}})\times j^*(\sisi{N_{Z/X}}{N_{X/Z}})\to \sisi{N_{T/X}}{N_{X/T}}\) es sobreyectivo, y que su núcleo está formado por los elementos que tienen coordenadas nulas sobre los \((\overline g_j,0)\) y los \((0,\overline h_k)\), y coordenadas iguales para todo \(i\) sobre \((\overline f_i,0)\) y \((0,\overline f_i)\). Por consiguiente, la sucesión es exacta también en el medio. Ahora bien, los tres \(\cO_T\)-módulos que intervienen son localmente libres de tipo finito; como \(\Tor_1^{\cO_X}(\cO_Y,\cO_Z)\) es de rango \(p\) por la demostración de i), se ve inmediatamente que la sucesión es exacta.

\subsection*{2.6}
Supongamos ahora que \(X\) es un esquema que posee un haz amplio. Todo subesquema cerrado \(Y\) de \(X\) posee entonces un haz amplio. Se sabe que en este caso «~existen resoluciones globales~», y que los \(K^\circ\) formados respectivamente con los haces localmente libres de tipo finito y con los complejos perfectos son idénticos (II, 2.2.9). Si, además, la inmersión \(i:Y\to X\) es regular, es un morfismo perfecto, y se puede definir un homomorfismo aditivo \(i_*:K^\circ(Y)\longrightarrow K^\circ(X)\) (IV 2.12), asociando a la clase de un complejo perfecto \(L^\bullet\) la clase del complejo \(\mathbb Ri_*(L^\bullet)=i_*(L^\bullet)\) (siendo \(i_*\) un funtor exacto).

\textbf{Proposición 2.7.} Sean \(X\) un esquema que posee un haz amplio, \(i:Y\to X\) una inmersión cerrada regular, y \(N\) el haz conormal de \(Y\) en \(X\). Entonces, para todo \(x\in K^\circ(Y)\), se tiene la relación
\begin{equation}\tag{2.7.1}
i^*(i_*(x))=x.\lambda_{-1}(N),
\end{equation}
donde \(N\) es la clase de \(N\) en \(K^\circ(Y)\), y \(\lambda_{-1}(N)=\sum_{i\ge0}(-1)^i\lambda^i(N)\).

\emph{Demostración.} Como los dos miembros de la relación que hay que demostrar son aditivos, podemos suponer que \(x\) es la clase de un \(\cO_Y\)-módulo localmente libre de tipo finito \(M\). El elemento \(i^*(i_*(x))\) es entonces la clase del complejo \(\mathbb L i^*(\mathbb Ri_*(M))\), obtenido tomando una resolución proyectiva del complejo \(\mathbb Ri_*(M)=i_*(M)\) y tensorizándola por \(\cO_Y\). Los haces de homología de este complejo son, pues, los \(\Tor_i^{\cO_X}(M,\cO_Y)\). Por 2.4, se tiene el isomorfismo
\[
\Tor_i^{\cO_X}(M,\cO_Y)
\xrightarrow{\sim}
M\otimes_{\cO_Y}\Tor_i^{\cO_X}(\cO_Y,\cO_Y).
\]
Por 2.5, se tiene el isomorfismo
\[
\Tor_i^{\cO_X}(\cO_Y,\cO_Y)
\xrightarrow{\sim}
\bigwedge^i N.
\]
Los \(\Tor_i^{\cO_X}(M,\cO_Y)\) son, por tanto, \(\cO_Y\)-módulos localmente libres de tipo finito, y la clase en \(K^\circ(Y)\) del complejo \(\mathbb L i^*(\mathbb Ri_*(M))\) es la suma alternada de las clases de sus haces de homología, es decir,
\[
i^*(i_*(x))=\sum_{i\ge0}(-1)^i
\operatorname{cl}\bigl(M\otimes_{\cO_Y}\underline{\Tor}_i^{\cO_X}(\cO_Y,\cO_Y)\bigr)
=\sum_{i\ge0}(-1)^i x.\lambda^i(N).
\]

\textbf{Corolario 2.8.} Bajo las hipótesis de 2.7, sean \(x,y\in K^\circ(Y)\); se tiene la relación
\begin{equation}\notag
i_*(xy\lambda_{-1}(N))=i_*(x).i_*(y).
\end{equation}
Por 2.7, \(y\lambda_{-1}(N)=i^*(i_*(y))\). Aplicando la fórmula de proyección, se obtiene
\[
i_*(xy\lambda_{-1}(N))
=i_*(x.i^*(i_*(y)))
=i_*(x).i_*(y).
\]
